\section{System Architecture and Design}
\label{sec:design}

Our system architecture bridges the gap between structured scientific records and raw consensus layer blobs through a deterministic four-stage pipeline: (1) Canonicalization, (2) Compression, (3) Field Modulus Packing, and (4) Three-Tiered Commitment Binding, integrated with a Dual-Horizon Archival Bridge.

\subsection{Canonical Serialization and Framing}
Reproducibility requires that identical scientific records always produce identical cryptographic commitments regardless of the runtime environment, CPU architecture, or JSON dictionary ordering. We enforce strict deterministic serialization adhering to RFC~8785 (JSON Canonicalization Scheme, JCS)~\cite{rfc8785}:
\begin{enumerate}
    \item Top-level keys are sorted lexicographically by UTF-16 code units.
    \item Whitespace between tokens is strictly eliminated (\texttt{separators=(',', ':')}).
    \item IEEE 754 floating-point numbers are normalized according to ECMAScript standard representations.
\end{enumerate}
Records are assembled into a contiguous record stream framed by 4-byte big-endian length prefixes:
\begin{equation}
S = \bigoplus_{i=1}^N \left( \text{len}(R_i)_{\text{uint32}} \parallel R_i \right)
\end{equation}

\subsection{Compression and Framing Strategy}
Structured scientific data (e.g., repeating keys in JSONL traces, tabular telemetry) exhibits high syntactic redundancy. We apply Zstandard (RFC~8878)~\cite{rfc8878} with configurable compression levels. Compression occurs over the framed record stream prior to field packing, yielding a compressed payload $C(S)$.

\subsection{Modulus-Safe Field Element Packing (\texttt{FE31})}
Because consensus clients reject field elements $\ge p$, we design a modulus-safe packing strategy denoted \texttt{FE31}. Each 128~KiB blob is partitioned into 4,096 field elements of 32 bytes each. In \texttt{FE31}, exactly 31 data bytes are placed in the lower 31 bytes of each field element, while the most significant byte (MSB) is zeroed (\texttt{0x00}):
\begin{equation}
\text{FE}_k = \texttt{0x00} \parallel D_{k,0} \parallel D_{k,1} \parallel \dots \parallel D_{k,30}
\end{equation}
Since $\text{MSB} = 0$, every resulting 32-byte integer satisfies $\text{FE}_k \le 2^{248} - 1 < p$, mathematically guaranteeing that invalid field element errors are impossible. Under \texttt{FE31}, each blob provides:
\begin{equation}
\text{UsableBytes}_{\text{FE31}} = 4,096 \times 31 = 126,976 \text{ bytes}
\end{equation}
Field Element 0 is reserved for a 31-byte domain-separated header (\texttt{magic="RBDA"}, version, codec, layout, blob index, blob count, chunk length, and content hash tag), leaving $126,945$ usable body bytes per blob.

\subsection{Three-Tiered Hybrid Commitments}
To support selective verification, whole-dataset integrity, and on-chain dispute games, we implement a three-tiered commitment hierarchy:
\begin{enumerate}
    \item \textbf{Tier 1: Blob Versioned Hashes ($H_v$)}: Generated by consensus layer KZG commitments and retrieved on-chain via the \texttt{BLOBHASH} opcode, proving that the exact published bytes were attested by Ethereum validators.
    \item \textbf{Tier 2: Canonical Content Hash ($H_c$)}: The Keccak-256 digest of the uncompressed canonical stream ($H_c = \text{keccak256}(S)$), providing an invariant whole-dataset fingerprint independent of compression codec or packing strategy.
    \item \textbf{Tier 3: Record-Level Merkle Tree Root ($R_m$)}: A position-binding binary Merkle tree constructed over individual canonical records. Sibling nodes are hashed with domain-separated prefixes:
    \begin{align}
    \text{Leaf}_i &= \text{keccak256}(\texttt{0x00} \parallel \text{keccak256}(R_i)) \\
    \text{Node} &= \text{keccak256}(\texttt{0x01} \parallel \text{Left} \parallel \text{Right})
    \end{align}
    This enables $O(\log N)$ inclusion proofs for selective record retrieval and on-chain fraud refutations.
\end{enumerate}

\subsection{Dual-Horizon Archival Bridge}
To address the 4,096-epoch consensus pruning horizon, the architecture decouples temporary publication from durable preservation:
\begin{itemize}
    \item At publication time, the publisher uploads the exact canonical payload to an external content-addressed store (e.g., IPFS or institutional CAS) to obtain a CIDv1 multihash.
    \item The publisher registers the 32-byte CID hash in \texttt{ResearchArchiveRegistry.sol}, linking it directly to the finalized Ethereum dataset record.
    \item Observers verify the dataset from blobs during the 18.2-day window; after blob pruning, observers retrieve the payload from the content-addressed store and verify it against the Ethereum on-chain $H_c$ and $R_m$ commitments.
\end{itemize}
